% GENERATED, edit the French unit or the translation table
\mxunit{automatismes-ch06}{Algebra, equations and inequalities}

\begin{fluencysheet}{Thu~14~Jan.~2027}{Integer powers (positive and negative)}
\mxfq{Warm-up · Quick}{PUI}{Calculate $2^{-3}$, $10^{-2}$ and $(-3)^2$.}{$\dfrac18$; $\num{0,01}$; $9$.}{}
\mxfq{Core · Method}{PUI}{Simplify $\dfrac{2^3\times2^{-5}}{2^{-1}}$.}{$\dfrac{2^{-2}}{2^{-1}}=2^{-2+1}=2^{-1}=\dfrac12$.}{}
\mxfq{Core + · Method}{PUI}{For $x\neq0$, write $(x^2)^{-3}\times x^7$ in the form $x^n$.}{$x^{-6}\times x^7=x^{1}=x$.}{}
\mxfq{Review · Quick}{STA}{Median of the series $4,\,7,\,9,\,10,\,12,\,15$?}{Even number of values: $\dfrac{9+10}{2}=\num{9,5}$.}{}
\mxfq{Thread · Identify}{SOL}{Which solid has a disc-shaped base, a curved lateral surface and a single vertex?}{A right circular cone.}{}
\end{fluencysheet}

\begin{fluencysheet}{Mon~18~Jan.~2027}{Square roots; $\sqrt{a^2}=|a|$}
\mxfq{Warm-up · Quick}{CAL}{Calculate $\sqrt{36}$ and $\sqrt{\num{0,09}}$.}{$6$; $\num{0,3}$.}{}
\mxfq{Core · Method}{CAL}{Write $\sqrt{72}$ in the form $a\sqrt2$ and calculate $\sqrt{(-5)^2}$.}{$\sqrt{72}=\sqrt{36\times2}=6\sqrt2$; $\sqrt{(-5)^2}=|{-}5|=5$.}{}
\mxfq{Core + · True/false}{CAL}{True or false? Justify. $\sqrt{9+16}=\sqrt9+\sqrt{16}$.}{False. $\sqrt{25}=5$ but $3+4=7$. Only the product ``splits'', for $a\geqslant0$ and $b\geqslant0$: $\sqrt{ab}=\sqrt a\sqrt b$.}{}
\mxfq{Review · Quick}{PUI}{Calculate $(3^2)^3$.}{$3^{2\times3}=3^6=729$.}{}
\mxfq{Thread · True/false}{DEM}{True or false? Justify. The diagonals of a rhombus are perpendicular.}{True. Property of the rhombus (its diagonals are also perpendicular bisectors of each other).}{}
\end{fluencysheet}

\begin{fluencysheet}{Tue~19~Jan.~2027}{Fractional expressions; expanded, reduced or factorised forms}
\mxfq{Warm-up · Quick}{CAL}{Calculate $\dfrac12+\dfrac13$.}{$\dfrac56$.}{}
\mxfq{Core · Method}{LIT}{Write as a single fraction $\dfrac1x+\dfrac1{x+1}$ (for $x\neq0$ and $x\neq-1$).}{$\dfrac{x+1+x}{x(x+1)}=\dfrac{2x+1}{x(x+1)}$.}{}
\mxfq{Core + · Method}{DEV}{Factorise $x^2-4$ then simplify $\dfrac{x^2-4}{x-2}$ for $x\neq2$.}{$x^2-4=(x-2)(x+2)$; so $\dfrac{x^2-4}{x-2}=x+2$.}{The choice of form depends on the goal: factorised to simplify, expanded to compare.}
\mxfq{Review · Quick}{SUB}{Calculate $2a-3b$ for $a=2$ and $b=-1$.}{$4+3=7$.}{}
\mxfq{Thread · Quick}{REP}{$ABCDEFGH$ is a cuboid with $AB=4$, $AD=3$, $AE=2$. In the frame $\pt{A;B,D,E}$, what are the coordinates of $G$?}{$G\pt{1;1;1}$: $G$ is the vertex opposite $A$ and $\overrightarrow{AG}=\overrightarrow{AB}+\overrightarrow{AD}+\overrightarrow{AE}$.}{A third coordinate appears in space. In this frame, the lengths $4$, $3$, $2$ do not appear in the coordinates.}
\end{fluencysheet}

\begin{fluencysheet}{Thu~21~Jan.~2027}{Inequalities: sum, multiplication by a real number}
\mxfq{Warm-up · Quick}{CMP}{Complete with $<$ or $>$: $-3\ \ldots\ -2$; $\dfrac13\ \ldots\ \num{0,3}$.}{$-3<-2$; $\dfrac13>\num{0,3}$ ($\num{0,333}\ldots>\num{0,3}$).}{}
\mxfq{Core · Method}{EQ}{We have $a<b$. Compare $-2a$ and $-2b$, then $a+5$ and $b+5$.}{$-2a>-2b$ (multiplying by a negative number reverses the inequality); $a+5<b+5$.}{}
\mxfq{Core + · Method}{EQ}{Given $2<x<3$ and $1<y<4$, find bounds for $x+y$ then $x-y$.}{$3<x+y<7$; for $x-y$: $-y\in\intoo{-4}{-1}$ so $-2<x-y<2$.}{Pitfall: subtracting the inequalities term by term.}
\mxfq{Review · Quick}{CAL}{Calculate $(-3)\times(-4)+2\times(-5)$.}{$12-10=2$.}{}
\mxfq{Thread · Problem}{DIV}{A $120$ cm by $84$ cm sheet is cut into identical squares, with no waste. What is the largest possible side length?}{$\pgcd\pt{120;84}=12$: squares of side $12$ cm ($10\times7=70$ squares).}{}
\end{fluencysheet}

\begin{fluencysheet}{Mon~25~Jan.~2027}{Comparing by difference and by quotient}
\mxfq{Warm-up · Quick}{CMP}{Compare $\dfrac23$ and $\dfrac34$ by studying the sign of the difference.}{$\dfrac23-\dfrac34=-\dfrac1{12}<0$: $\dfrac23<\dfrac34$.}{}
\mxfq{Core · Method}{LIT}{Compare $x^2+1$ and $2x$ for every real number $x$.}{$x^2+1-2x=(x-1)^2\geqslant0$: $x^2+1\geqslant2x$, with equality if and only if $x=1$.}{}
\mxfq{Core + · Method}{CMP}{Compare $\dfrac75$ and $\dfrac86$ by quotient (positive numbers).}{$\dfrac{7/5}{8/6}=\dfrac{42}{40}>1$, so $\dfrac75>\dfrac86$.}{}
\mxfq{Review · Table}{LEC}{Monthly rainfall (fictional data).\par\smallskip\mxarraystretch{1.15}\begin{tabular}{|c|c|c|c|c|}\hline
Month & Jan. & Feb. & Mar. & Apr. \\ \hline
Rain (mm) & 280 & 300 & 240 & 120 \\ \hline
\end{tabular}\par Which month is the wettest? What is the total over these $4$ months?}{February ($300$ mm); $940$ mm.}{}
\mxfq{Thread · Table → formula}{REG}{$x$: $-2,\,-1,\,0,\,1$ and $f(x)$: $4,\,1,\,0,\,1$. Is the function affine? Suggest a formula.}{No (the changes are not constant). $f(x)=x^2$ works.}{Test: $f(1)-f(0)=1$ but $f(0)-f(-1)=-1$.}
\end{fluencysheet}

\begin{fluencysheet}{Tue~26~Jan.~2027}{First-degree equations and inequalities; $x^2=a$; zero product}
\mxfq{Warm-up · Quick}{EQ}{Solve $3x-5=7$.}{$x=4$.}{}
\mxfq{Core · Method}{EQ}{Solve $2x+3\leqslant7-x$.}{$3x\leqslant4$: $x\leqslant\dfrac43$; $S=\left]-\infty;\dfrac43\right]$.}{}
\mxfq{Core + · Method}{EQ}{Solve $(x-2)(2x+6)=0$ then $x^2=7$.}{Zero product: $x=2$ or $x=-3$; $x^2=7$: $x=\sqrt7$ or $x=-\sqrt7$.}{}
\mxfq{Review · Quick}{CMP}{Compare $\num{1,4}$ and $\sqrt2$ without a calculator.}{$\num{1,4}^2=\num{1,96}<2$ so $\num{1,4}<\sqrt2$ (the square function is increasing on $\intfo{0}{+\infty}$).}{}
\mxfq{Thread · Quick}{GEO}{$ABC$ is isosceles at $A$ with $\widehat{BAC}=40\degre$. What is the measure of $\widehat{ABC}$?}{$(180-40)\div2=70\degre$.}{}
\end{fluencysheet}

\begin{fluencysheet}{Thu~28~Jan.~2027}{Quotient equation $\frac{A(x)}{B(x)}=k$}
\mxfq{Warm-up · Quick}{EQ}{For which value of $x$ is the expression $\dfrac{1}{x-3}$ undefined?}{$x=3$.}{}
\mxfq{Core · Method}{EQ}{Solve $\dfrac{2x-1}{x+3}=0$.}{$x\neq-3$ and $2x-1=0$: $x=\dfrac12$.}{}
\mxfq{Core + · Method}{EQ}{Solve $\dfrac{x+1}{x-2}=3$.}{$x\neq2$; $x+1=3(x-2)$, that is $x=\dfrac72$. Check: $\dfrac{\num{4,5}}{\num{1,5}}=3$.}{}
\mxfq{Review · Expand}{DEV}{Expand and simplify $(x+2)(x-5)$.}{$x^2-3x-10$.}{}
\mxfq{Thread · Code reading}{ALG}{What do $g(3)$ and $g(-5)$ return?\mxpseudo{def g(x):\\ \mxtab if x \textgreater{} 0:\\ \mxtab\mxtab return 2*x\\ \mxtab else:\\ \mxtab\mxtab return -x}}{$g(3)=6$ and $g(-5)=5$.}{The condition is evaluated before the calculation: the branch depends on the sign of $x$.}
\end{fluencysheet}

\begin{fluencysheet}{Mon~1\textsuperscript{st}~Feb.~2027}{Sign of a product or quotient; sign tables}
\mxfq{Warm-up · Quick}{EQ}{Sign of $2x-4$?}{Negative for $x<2$, zero at $2$, positive for $x>2$.}{}
\mxfq{Core · Method}{EQ}{Draw the sign table of $(x-1)(x+3)$.}{Roots $-3$ and $1$: positive on $\intoo{-\infty}{-3}$, negative on $\intoo{-3}{1}$, positive on $\intoo{1}{+\infty}$.}{Sign of the product: rule of signs, column by column.}
\mxfq{Core + · Method}{EQ}{Solve $\dfrac{x+2}{x-1}\geqslant0$.}{Numerator positive for $x\geqslant-2$, denominator positive for $x>1$; $S=\intof{-\infty}{-2}\cup\intoo{1}{+\infty}$ ($1$ is excluded).}{}
\mxfq{Review · Quick}{PRB}{Pack of $32$ cards: probability of drawing a heart? a king?}{$\dfrac{8}{32}=\dfrac14$; $\dfrac{4}{32}=\dfrac18$.}{}
\mxfq{Thread · Disprove}{LIT}{Disprove: ``for every real number $x$, $x^2\geqslant x$''.}{Counterexample: $x=\dfrac12$: $x^2=\dfrac14<\dfrac12=x$.}{A single counterexample is enough to disprove a universal statement.}
\end{fluencysheet}
